\section{Optimal Transport for Distributional Selection} \label{sec:28-background-ot}

\S\ref{sec:27-background-influence} corrected the per-sample weights but left property~(2) open: even with accurate weights, ranking and selecting independently cannot guarantee that the chosen subset covers the weighted distribution.
This section develops optimal transport as the tool that closes that gap, formalizing selection as a distributional problem rather than a ranking one.

\begin{foldable}
  Optimal transport has been widely applied across \ac{NLP} and machine learning, from Word Mover's Distance~\cite{kusner2015word} and Wasserstein autoencoders~\cite{tolstikhin2018wasserstein} to domain adaptation~\cite{courty2016optimal} and Wasserstein \acp{GAN}~\cite{arjovsky2017wasserstein}, before we turn to its role for distributional coverage in dataset distillation.
\end{foldable}

\begin{foldable}
  \textbf{The distributional coverage problem.}
  Property~(2) requires that the selected subset jointly represent the importance-weighted distribution, not just consist of high-importance examples.
  At extreme compression, when $K \ll N$, a natural objective emerges: select a compact subset that minimises the distance between the reweighted source distribution and the selected subset.
  This is a \textit{distributional} problem: it cannot be solved by ranking examples independently.
\end{foldable}

\begin{foldable}
  \textbf{Why independent selection fails architecturally.}
  Any rank-and-select strategy decomposes as a sum of independent per-sample decisions, invariant to relationships between selected examples.
  The selection of one example does not constrain or inform the selection of another; the algorithm ignores pairwise distances and redundancy in feature space.
  By contrast, a distributional objective penalizes redundancy and rewards coverage: if two selected examples are similar and both have high importance, a distributional selector can free capacity by deselecting one and using that slot for a distinct uncovered region.
  Architecturally, any greedy or clustering-based method that optimizes per-sample scores independently cannot solve a joint distributional problem.
\end{foldable}

\begin{foldable}
  \textbf{Alternative divergences and their limits.}
  Before naming a specific divergence, two natural candidates from the wider literature are worth ruling out for the regime here.
  Maximum mean discrepancy~\cite{gretton2012kernel} measures distance between distributions by comparing their mean embeddings under a kernel, but the comparison is symmetric and aggregate.
  It does not produce a sample-level assignment from source to surrogate, so selection on top of it requires an auxiliary heuristic that re-introduces the independent-selection problem of the previous foldable.
  Kullback--Leibler and Jensen--Shannon divergences require estimating the densities of source and surrogate, which is unreliable in the low-data regime~\cite{paninski2003estimation}; on top of that they become infinite when the supports are disjoint, which is the rule rather than the exception when $K \ll N$.
\end{foldable}

\begin{foldable}
  \textbf{Wasserstein distance and optimal transport.}
  Optimal transport formalizes the distributional coverage problem geometrically as a transportation problem.
  Imagine mass is distributed over the training set according to importance weights $w_n$; we want to transport this mass to a sparse surrogate of size $K$.
  The cost of moving one unit of mass from source example $z_n$ to surrogate example $\tilde{z}_c$ is the feature-space distance $d(z_n, \tilde{z}_c)$ (\eg, Euclidean distance in embedding space).
  The optimal transport problem finds a coupling (transport plan) $\gamma$ that moves all importance mass from source to surrogate at minimum total cost.
  The Wasserstein-$p$ distance~\cite{villani2008optimal} quantifies this minimum cost:
  \begin{equation}
    \label{eq:wasserstein}
    W_p(\mu, \nu) = \left(\inf_{\gamma \in \Gamma(\mu,\nu)} \int d(z, z')^p \, \text{d}\gamma(z, z') \right)^{1/p},
  \end{equation}
  where $\Gamma(\mu, \nu)$ is the set of all couplings (joint distributions) with marginals $\mu$ (source) and $\nu$ (surrogate).
  Intuitively, a small Wasserstein distance means the surrogate can cover the importance-weighted source cheaply: few examples have to travel far.
  Geometrically, $W_p(\mu, \nu)$ is the minimum cost of transporting probability mass from $\mu$ to $\nu$.
  Where \ac{KL} and Jensen-Shannon blow up under disjoint supports, Wasserstein keeps producing a meaningful distance, which matters once $K \ll N$.
  This makes Wasserstein the natural choice for measuring distributional coverage in dataset distillation.
\end{foldable}

\begin{foldable}
  \textbf{Discrete optimal transport formulation.}
  In dataset distillation, both source and target are discrete, finite sets and can be formulated as an \ac{OT} problem.
  Let the source be $\{(z_n, w_n)\}_{n=1}^{N}$ where $w_n$ are normalized importance weights (each $w_n \geq 0$, $\sum_n w_n = 1$).
  The candidate pool (synthetic examples or a large real example pool) is $\{(\tilde{z}_c, 1/C)\}_{c=1}^{C}$ with uniform weights (each candidate has equal potential budget $1/C$).
  The discrete optimal transport problem is a distributional selection problem: given weighted source examples and an unweighted candidate pool, find a transport plan that moves all source mass to the candidate pool at minimum cost.
  Formally, seek a transport plan $\gamma^* \in \mathbb{R}^{N \times C}_{\geq 0}$ minimising total transport cost subject to marginal constraints:
  \begin{equation}
    \gamma^* = \arg\min_{\gamma \geq 0} \langle C_{\text{dist}}, \gamma \rangle
    \quad \text{s.t.} \quad \gamma \mathbf{1}_C = \mathbf{w}, \quad
    \gamma^\top \mathbf{1}_N = \tfrac{1}{C}\mathbf{1}_C,
    \label{eq:discrete-ot}
  \end{equation}
  where $C_{\text{dist},nc} = d(z_n, \tilde{z}_c)^2$ is the squared cost of transporting mass from source example $z_n$ to candidate example $\tilde{z}_c$ (typically Euclidean distance in embedding space).
  The constraints enforce the transportation semantics: all importance mass leaves the source ($\gamma \mathbf{1}_C = \mathbf{w}$) and equal weight arrives at each candidate ($\gamma^\top \mathbf{1}_N = \tfrac{1}{C}$).
  The entry $\gamma^*_{nc}$ is a fractional transport weight encoding how much importance mass from source $z_n$ is optimally routed to candidate $\tilde{z}_c$.
  After solving, the $K$ candidates with the largest total assigned mass $m_c = \sum_n \gamma^*_{nc}$ are selected as the surrogate.
  Finally, \ac{OT} selects jointly to minimise distributional divergence, not independently to maximise per-example scores.
\end{foldable}

\begin{foldable}
  \textbf{Sinkhorn algorithm for tractable computation.}
  Solving discrete \ac{OT} exactly via linear programming is expensive for large pools (quadratic or worse in $N + C$).
  The Sinkhorn algorithm~\cite{cuturi2013sinkhorn} reformulates the problem with entropic regularization: adding a penalty term $\varepsilon \, \text{KL}(\gamma \| \mathbf{w} \otimes \tfrac{1}{C}\mathbf{1}_C)$ to the objective.
  This transforms the problem from a hard-constrained LP into a strongly convex one with a simple closed-form update: alternating row and column normalization.
  The update is simple: normalize rows of $\gamma$ to satisfy the source marginal, then normalize columns to satisfy the target marginal, repeat until convergence.
  Each iteration costs $O(N \cdot C)$, and convergence occurs in tens of iterations, making the algorithm practical for large pools.
  The regularised plan $\gamma_\varepsilon^*$ converges to the exact unregularised \ac{OT} optimum $\gamma^*$ as $\varepsilon \to 0$, providing a principled approximation-accuracy tradeoff.
  The regularization parameter $\varepsilon$ is a tuning knob: large $\varepsilon$ spreads mass diffusely (smooth, entropy-maximizing plan), small $\varepsilon$ produces concentrated plans closer to exact \ac{OT}.
\end{foldable}

\begin{foldable}
  \textbf{Concrete example: \ac{OT} vs. greedy and clustering.}
  Suppose a training set has six examples in two semantic regions: Region~A has four examples each with importance $0.25$ (total $1.0$); Region~B has two examples each with importance $0.05$ (total $0.10$).
  From a candidate pool of four examples (two near each region), select $K=2$ for the surrogate.
\end{foldable}

\begin{foldable}
  \textbf{$k$-means clustering} partitions candidates by geometric proximity without consulting importance weights.
  If Region~B candidates are geometrically cohesive, they form one cluster; Region~A candidates form the other.
  $k$-means selects one centroid per cluster: one Region~A candidate and one Region~B candidate.
  The surrogate allocates $50\%$ capacity to Region~B despite Region~B carrying only $9\%$ of the importance mass, a factor-of-five mismatch.
\end{foldable}

\begin{foldable}
  \textbf{Greedy $k$-centre selection} minimises worst-case distance: it chooses the two geometric extremes, again one Region~A and one Region~B candidate.
  Same outcome: both methods conflate geometric coverage with distributional fidelity, and when the two decouple, as in text where low-importance examples cluster in distinctive surface-form niches, geometry actively misleads.
\end{foldable}

\begin{foldable}
  \textbf{Optimal transport} with weights $\mathbf{w} = (0.25, 0.25, 0.25, 0.25, 0.05, 0.05)$ constructs a plan that moves Region~A's mass (1.0 units) and Region~B's mass (0.10 units) to the candidate pool.
  The Sinkhorn plan assigns $\approx 0.50$ units to each Region~A candidate and $\approx 0.05$ units to each Region~B candidate.
  Selecting the two candidates with largest total assigned mass returns both Region~A candidates and neither Region~B candidate.
  The surrogate omits Region~B entirely, exactly the outcome that minimises Wasserstein distance from the importance-reweighted source.
  The advantage of \ac{OT} is not geometric coverage but \textit{importance-weighted} coverage: it concentrates surrogate capacity on the regions that matter most to the learning objective.
\end{foldable}
